\subsection{The height of a valuation}

Let G be a totally ordered group. Given two isolated subgroups H and H' of G, one of them is contained in the order: for otherwise there would exist a positive element x of H not belonging to H' and a positive element \(x'\) of H' not belonging to H; let, for example \(x \geqslant x'\) ; as H is isolated, \(x' \in H\) , which is a contradiction.

This also follows from Proposition 4 of no. 3 and the Corollary to Proposition 1 of no. 1, taking account of the fact that every totally ordered group is the order group of a valuation (§ 3, no. 4, Example 6).

DEFINITION 2. Let G be a totally ordered group. If the number of isolated subgroups of G distinct from G is finite and equal to n, G is said to be of height n.~If this number is infinite, G is said to be of infinite height.

\textbf{Examples}\par

\begin{enumerate}
\def\labelenumi{(\arabic{enumi})}
\item
  The height of the group \(\mathbf{G} = \{0\}\) is 0.
\item
  The groups \(\mathbf{Z}\) and \(\mathbf{R}\) are of height 1.
\item
  Let G be a totally ordered group and H an isolated subgroup of G. If \(h(\mathrm{H})\) and \(h(\mathrm{G}/\mathrm{H})\) denote the heights of the totally ordered groups H and G/H, then
\end{enumerate}

\[
h (\mathbf {G}) = h (\mathbf {H}) + h (\mathbf {G} / \mathbf {H}),\tag{2}
\]

since the set of isolated subgroups of G is totally ordered by inclusion. In particular, if G is the lexicographic product of two totally ordered groups H and H', then

\[
h (\mathbf {G}) = h (\mathbf {H}) + h \left(\mathbf {H} ^ {\prime}\right)\tag{3}
\]

(cf.~no. 2, Example 2); thus the lexicographic product \(\mathbf{Z} \times \mathbf{Z}\) is of height 2.

On the other hand the height of \(\mathbf{Z} \times \mathbf{Z}\) ordered by embedding in \(\mathbf{R}\) (cf.~§ 3, no. 4, end of Example 6) is equal to 1 (cf.~Proposition 8 below).

DEFINITION 3. The height of the order group of a valuation is called the height of that valuation.

For example a discrete valuation is of height 1. Only improper valuations are of height 0. Propositions 1 and 4 imply:

PROPOSITION 5. The height of a valuation is equal to the number of non-zero prime ideals in its ring.

